\documentclass[12pt]{article}
\usepackage{geometry}
\geometry{margin=1in}
\usepackage{float}
\usepackage{graphicx}
\usepackage{hyperref}

\title{Jaw Models Explanation}
\author{Zebrafish Osteoarthritis Project}
\date{\today}

\begin{document}

\maketitle

\section{Introduction}
One of the best ways to image bones in an animal (here, zebrafish) is with a CT scanner.
This is non-invasive and provides us with detailed 3D images that we can do all sorts of things with.

The CT scan is a greyscale (bit depth of 1) image.
Normally, investigating the properties of a bone requires us to:
\begin{itemize}
    \item Load the CT scan image into some imaging program
    \item Look through each slice of the scan, manually labelling each voxel as bone or not
    \item Export these labels as a segmentation mask
    \item Load the mask into some other program and perform analysis on it
\end{itemize}
The CT scan is a large image (it will have a dimension of around $2000\times 500\times 500$ voxels),
so making the segmentation is time-consuming and can be error-prone (humans get tired from looking at a
screen for hours).
It is also not fully reproducible --- different humans will have different opinions on e.g.
whether one voxel represents bone or not, where one bone ends and another begins, and where things like
the edges of bones lie.

To solve some of these problems, AI methods can be used to segment the bones of interest out from
a CT scan.

\section{Analysis Pipeline}
This is what happens to a CT scan when we put it through our analysis pipeline:
\begin{figure}[H]
    \includegraphics[width=\textwidth]{assets/fish_procedure.png}
    \label{analysis pipeline}
    \caption{
        The analysis pipeline.
        A CT image is fed into the jaw location model, which finds the rough location of the jaw.
        The scan is then cropped around this location, and the cropped scan fed into the segmentation model.
        The segmentation model works by splitting the input image into patches (see Section.~\ref{Jaw Segmentation}),
        running inference on each patch and then aggregating them together.
        The final jaw segmentation mask is created by thresholding the model's output to get a binary image,
        and taking the largest connected component in the binary image.
        We can then do downstream analysis on this segmentation mask (meshing, FE analysis, ...)
    }

\end{figure}

\section{Jaw Location}
\label{Jaw Location}
We want to feed our CT scan into some model pipeline, and have it output a segmentation mask (or mesh) of the jawbone.

However, as mentioned above each image is large (around 1.6GB on disk in practice for our $2000\times 640\times 640$ images).
Training a model on these images would probably be
\begin{enumerate}
    \item inefficient (most of the CT scan is not the jaw, so the model would spend most of its time learning that 99.9\% of the scan is not jaw)
    \item slow or expensive (we would either have to pay to use a very large GPU that can hold tens to hundreds of GB training data in memory, or read the data from disk at every epoch which would be very slow)
\end{enumerate}
To deal with this, we trained the jaw segmentation model (Section \ref{Jaw Segmentation}) on cropped-out data -
we manually cropped a region around the jaw in both the mask and the training image, and fed these pairs into the model.

This has an unfortunate side effect --- it means that the jaw segmentation model (probably) won't work well if we feed it
entire CT scans.
We don't want to spend lots of time manually cropping out the jaws when we take new data from a beamline --- this would defeat
some of the point of our time-saving segmentation pipeline.
To get around this, we trained a model to locate the jaw automatically.

\subsection{How the Jaw Location Model Works}
\subsubsection{Downsampling Scan}
The CT scans are too large to fit in memory (as described above),
so we can't train the locator model on the full-sized, full-resolution scan.
However, the location model doesn't need to see all the fine detail of the fish - it just needs to see its general structure.
We can therefore downsample its training data, and still have a model that is capable of finding the right location in the scan.
\begin{figure}[H]
    \centering
    \label{downsampling}
    \includegraphics[width=0.48\textwidth]{../../script_output/jaw_location/example_locator/test_centroid_truth.png}
    \includegraphics[width=0.48\textwidth]{../../script_output/jaw_location/example_locator/test_centroid_downsampled.png}
    \caption{
        The original (left) and downsampled (right) images.
        The downsampled image is approximately a sixteenth the size of the original image, but still contains
        enough context for the model to locate the jaw.
        The location of the jaw centroid is marked in red.
    }
\end{figure}
Once the jaw is located in the downsampled scan, it is a simple matter to convert the jaw location co-ordinate in the downsampled scan
to the corresponding co-ordinate in the full-resolution scan. E.g. if we downsample the image by 4 times (in each direction), we just
multiply the location co-ords by 4.

\subsubsection {Object Location}
The intuitive way to locate an object in an image would be to train a model to learn the $\left(x, y, z\right)$ co-ordinates of the object
from the image.
In practice, this doesn't work very well - it can be hard to train a model to learn the co-ordinates without overfitting
(part of this is that the output won't have the ``inductive biases'' of a convolutional network - we lose things like translational symmetry
when our model learns co-ords from an image).
Instead, we actually train a U-Net (which is a segmentation model) for the jaw location.
This works by placing a Gaussian heatmap at the centre of the jaw--- the model is tasked to segment this Gaussian out
from the CT scan.
As well as being easier to train, additional benefits of the heatmap are:
\begin{itemize}
    \item Heatmaps easily generalise to identifying multiple objects in one image
    \item We get a free uncertainty measure of the location from the heatmap
    \item We can derive things like bounding boxes easily from heatmaps
    \item Heatmaps can handle overlapping objects (if we e.g. wanted to, later, identify )
    \item We can get sub-pixel precision, if our heatmap is $<1$px in size
\end{itemize}
During training, the heatmap is gradually shrunk to give us better precision.
We start with a large heatmap, since that is easy for the model to learn.
Once it is performing well enough, we shrink it (down to some minimum size).

The segmentation model (Section~\ref{Jaw Segmentation}) needs to be good to give meaningful results.
In contrast, the location model only needs to be ``good enough'' that we can crop the full jaw from its
output--- it doesn't matter if we're a few pixels off, as long as we still get the jaw in the cropped output.
\begin{figure}[H]
    \centering
    \includegraphics[width=0.45\textwidth]{../../script\_output/jaw\_location/example\_locator/heatmap\_epoch\_49\_sigma\_7.200.png}
    \includegraphics[width=0.45\textwidth]{../../script\_output/jaw\_location/example\_locator/heatmap\_epoch\_173\_sigma\_3.444.png}
    \caption{
        The jaw location model is a segmentation model - we learn to segment out a Gaussian heatmap placed at the jaw's centroid.
        The heatmap gradually shrinks as the model learns.
    }
\end{figure}

\begin{figure}[H]
    \centering
    \includegraphics[width=\textwidth]{../../script\_output/jaw\_location/example\_locator/losses.png}
    \caption{
        The loss function can look funny---
        it decreases as we learn to find the large heatmap,
        then there are bumps as the heatmap shrinks and the model must re-learn.
        The learning rate in this example was a little too high, but you can sort of see the bumps
        in the training loss as the heatmap shrinks.
    }
\end{figure}

\section{Jaw Segmentation}
The segmentation model is fed a cropped scan, and outputs a matching segmentation mask.
As part of the inference pipeline, we threshold the model's output and clean it up so that the segmentation
mask only contains one object--- we don't necessarily have to do this, but it does give more sensible segmentation masks.
\label{Jaw Segmentation}
\subsection{Model Architecture}
Here's a simplified view of the model architecture:
\begin{figure}[H]
    \centering
    \includegraphics[width=0.8\textwidth]{assets/simple\_net\_diagram.png}
    \caption{
        A simplified view of what happens in a U-Net model.
        Patches are extracted from the input image (see~\ref{Patch-based Training}).
        The patch is fed through successive layers,
        each of which applies convolution operation(s).
        The convolution operations are basically sliding-window multiplications with a small image (the ``filter'').
        The network learns which convolution filters will give the correct output.
        This results in, roughly, the model learning to identify features at different length
        scales and how to translate them into a segmentation mask with the right shape at the right location.
    }
\end{figure}
A more detailed view, showing the actual operations:
\begin{figure}[H]
    \centering
    \includegraphics[width=0.8\textwidth]{assets/net_diagram.png}
    \caption{
        A more detailed view of what happens to the patch.
        The patch starts as a greyscale image (1 channel, 160 pixels long/wide/deep).
        Each layer starts with a convolution, to increase the number of channels (black arrow) and
        a max-pooling (green) which reduces the spatial dimension.
        The skip connections (dotted arrow) are modulated by an attention mechanism ($\sigma$),
        which uses the feature map from a deeper layer as the gating signal.
        Intuitively this means that the model uses the deeper (coarser, more semantically-aware) information
        to choose which information (finer, higher-resolution) is passed to the decoder in the skip connection.
    }
\end{figure}

\subsection{Patch-based Training}
\label{Patch-based Training}
A U-Net model has a fixed-sized input.
This is fine, as long as we know that its input will always be the same size.
However, in our case we may want to fine-tune the model to segment other bones, which might have a different input size
(e.g. the spine, which will probably be a very long and thin image).
To overcome this, we can split the input image into patches and then feed them into the model.
We lose some local context here, but it is ok as long as the patches are large enough to contain all the context needed for segmentation.
It can also overcome GPU memory limitations and force the model to focus more on local details, which might be desirable
(e.g. we might want to force the model to segment based on the surrounding tissue, not by looking for an object in the fish head).

The patches we used here are almost as big as the entire input image, so there wasn't a big memory saving.
However, the patch-based pipeline means that, in principle, might be able to feed the model an entire CT scan which
we could then run inference on. The model would ideally just segment the jawbone out from the scan, and leave everything else.
This doesn't work very well in practice though (the model hasn't seen much non-jaw fish, so it randomly segments out other bits).

We get the final segmentation mask by stitching together the model's output patches.
In the regions where the patches overlap, we just average the predictions from each patch (they tend to agree pretty well).

\subsection{Model Output}
The model outputs probabilities (sort of --- I think there's some technical reason why they aren't really probabilities, but it outputs a number between 0 and 1 for each voxel),
which we threshold to get the final prediction mask. We then clean it up - we know the jaw has to be a single, contiguous bone, so
we remove any voxels that aren't attached to the largest connected component in the segmentation mask.

Here are some slices through a jaw:
\begin{figure}[H]
    \includegraphics[width=0.45\textwidth]{../../script\_output/jaw\_models/example\_segmenter/train\_output/test\_pred.png}
    \includegraphics[width=0.45\textwidth]{../../script\_output/jaw\_models/example\_segmenter/train\_output/test\_truth.png}
    \caption{
        Jaw segmentation prediction vs truth; slices through the thing
    }
\end{figure}

\begin{figure}
    \centering
    \includegraphics[trim=100 100 140 0, clip, width=0.75\textwidth]{assets/wt.png}
    \includegraphics[trim=100 100 140 0, clip, width=0.75\textwidth]{assets/mutant.png}
    \includegraphics[trim=100 100 140 0, clip, width=0.75\textwidth]{assets/mutant3.png}
    \includegraphics[trim=100 100 140 0, clip, width=0.75\textwidth]{assets/mutant2.png}
    \label{jaws}
    \caption{
        Some examples of projections of the jaws.
        The colourbar shows greyscale value (density).
        These were extracted from Wahab's TIFS dataset on the RDSF, by putting the TIFS through
        our location/segmentation/thresholding pipeline.
        I think some of these are mutants.
    }
\end{figure}

\begin{figure}
    \includegraphics[width=\textwidth]{../../script\_output/jaw\_models/example\_segmenter/train\_output/loss.png}
    \caption{
        Jaw segmentation loss. Not sure what the big spikes in the validation loss are.
    }
\end{figure}

\section{Future Work}
There's a decent amount of stuff you could do with these models.
The obvious thing is to use them to segment out lots of jaws then do some useful analysis with the segmentation masks.

There are also a few obvious data science extensions, which I've briefly outlined below.
There's more detail on these, intended for a new person starting work on this project,
at \url{http://www.github.com/JGIBristol/Zebrafish_Osteoarthritis/issues}.
Others may be more immediately useful or interesting, depending on what's going on in the research group.

\subsection{Shape Analysis}
\subsubsection{Fine-tuning on Other Bones}
We want to fine-tune the model(s) to segment out other bones, e.g. the spine.
I've written a rough example showing this on another dataset; all we need to do is investigate
how well fine-tuning works on the new spine dataset.

\subsubsection{Auto Mesh}
At the moment we mesh the segmentations manually.
This is quite time-consuming; it should be possible to train a model to do this for us.
Cubic meshing is quite easy, once we have the segmentation mask.
A more complete workflow might take us straight from CT scan to a tetrahedral mesh.

\subsubsection{Shapeworks}
Software exists for automatic shape analysis.
We can do semilandmarking, normalisation and PCA to find out where the major shape differences
are located on the jaws.
This is useful for answering questions about how and where age/mutations/sex affect the bones.

\subsubsection{Multimodal Model for Multiple Bones}
The end goal would be a model that looks at the entire CT scan and segments out all the bones
with different labels. Perhaps there have been more recent advances in foundation models for
CT images which might help here.

\subsection{Density Analysis}
\subsubsection{Greyscale analysis}
We can see the muscle attachments in the images (Fig~\ref{jaws}).
They show up as regions of denser bone; we might want to investigate how the density at these
attachment sites changes, in absolute terms and relative to the rest of the bone, with age/mutation.

\subsubsection{Radiomic Feature Selection}
Radiomic features are things like greyscale patchiness/size/shape that are useful for radiologists.
We can extract radiomic features from our jaw dataset and see which ones are most useful for describing
the variation between aged/young/mutant fish.
This is useful since we can make hypotheses about the effect of density, and its distribution, and how
that affects fracture risk and bone optimality.

\subsubsection{Find Muscle Attachments}
You can see the muscle attachments in Fig~\ref{jaws}. To do the bone density analyis, we might
want to precisely locate them.
We could either do this by trying to come up with some simple greyscale-based rules,
or with another ML model.

\subsection{Existing Model Improvements}
\subsubsection{Model Ensemble}
At the moment, we just use one model for each task. There is some randomness involved in the training process, though
and we can exploit that by using an ensemble of models. This should, in some sense, average out the bias from our single model
and give us better results (e.g. if the segmentation model systematically under/overpredicts, an ensemble might help with this).
We can get this ``for free'', since we don't need to do anything special --- just train a load more models, then come up with
some way of combining their predictions (e.g. averaging, or some voting system).

\subsubsection{Ablation study}
As part of the fine-tuning study, I tried turning off the segmentation model's attention mechanism in selected layers.
The attention mechanisms suppress information passing though the skip connections --- when we turn these off, the model loses the
ability to ``focus on the important bits''.
Interestingly, it looked like there was very little effect from turning the attention in layers 1, 4 and 5 off.
However, turning the attention mechanism off in layers 2 or 3 had a huge effect:
\begin{figure}[H]
    \includegraphics[width=0.8\textwidth]{../../script\_output/ablation/example\_segmenter/ablation\_273\_2.png}
    \includegraphics[width=0.8\textwidth]{../../script\_output/ablation/example\_segmenter/ablation\_273\_3.png}
    \caption{
        Ablation study - turning layers 2 or 3 ``off'' seems to segment out all the bones in the fish head.
        This doesn't happen with the other layers, which is suggestive that the attention isn't doing much here.
        It is also suggestive that something is going on where the model is separately learning the shape and density
        of the target bone --- could we use this to understand how the model works, or to inform our fine-tuning studies?
    }
\end{figure}
NB I think one of STFC's priorities at the moment is on ``explainable AI'', so this might be a good thing to mention in grant applications.

\end{document}
